\chapter*{Environment Tests}
\section*{Result Environments}
\subsection*{Theorem Tests}
\begin{theorem}
    In any right triangle, the area of the square whose side is the hypotenuse is equal to the sum of the areas of the squares whose sides are the two legs.
    \[
        a^2 + b^2 = c^2
    \]
\end{theorem}

\begin{theorem}[Pythagorean Theorem]\label{thrm:test}
    In any right triangle, the area of the square whose side is the hypotenuse is equal to the sum of the areas of the squares whose sides are the two legs.
    \[
        a^2 + b^2 = c^2
    \]
\end{theorem}

See \myref{thrm:test}.

\subsection*{Lemma Tests}
\begin{lemma}
    Given two integers $n$ and $d \neq 0$, there exist unique integers $q$ and $r$ such that $n = qd + r$ and $0 \leq r < |d|$.
\end{lemma}

\begin{lemma}[B\'ezout]\label{lemma:test}
    Given two integers $n$ and $d \neq 0$, there exist unique integers $q$ and $r$ such that $n = qd + r$ and $0 \leq r < |d|$.
\end{lemma}

See \myref{lemma:test}.

\subsection*{Proposition Tests}
\begin{proposition}
    Let $m$ and $n$ be non-zero integers. Then
    \[
        |mn| = \mathrm{gcd}(m,n) \times \mathrm{lcm}(m,n).
    \]
\end{proposition}

\begin{proposition}[Relationship Between GCD and LCM]\label{prop:test}
    Let $m$ and $n$ be non-zero integers. Then
    \[
        |mn| = \mathrm{gcd}(m,n) \times \mathrm{lcm}(m,n).
    \]
\end{proposition}

See \myref{prop:test}.

\subsection*{Corollary Tests}
\begin{corollary}
    In any right triangle, the area of the square whose side is the hypotenuse is equal to the sum of the areas of the squares whose sides are the two legs.
    \[
        a^2 + b^2 = c^2
    \]
\end{corollary}

\begin{corollary}[Pythagorean Theorem]\label{coro:test}
    In any right triangle, the area of the square whose side is the hypotenuse is equal to the sum of the areas of the squares whose sides are the two legs.
    \[
        a^2 + b^2 = c^2
    \]
\end{corollary}

See \myref{coro:test}.

\subsection*{Definition Tests}
\begin{definition}
    An \term{inhomogeneous linear differential equation} has the form
    \[
        L[v] = f
    \]
    where $L$ is a linear differential operator, $v$ is the dependent
    variable, and $f$ is a given non-zero function of the independent
    variables alone.
\end{definition}

\begin{definition}[Inhomogeneous Linear DE]\label{def:test}
    An \term{inhomogeneous linear differential equation} has the form
    \[
        L[v] = f
    \]
    where $L$ is a linear differential operator, $v$ is the dependent
    variable, and $f$ is a given non-zero function of the independent
    variables alone.
\end{definition}

See \myref{def:test}.

\subsection*{Axiom Tests}
\begin{axiom}
    For any $x, y \in \R$ we have $x + y = y + x$.
\end{axiom}

\begin{axiom}[Commutivity of Addition]\label{axiom:test}
    For any $x, y \in \R$ we have $x + y = y + x$.
\end{axiom}

See \myref{axiom:test}.

\section*{Questions Environments}
\subsection*{Exercises Tests}
\begin{exercise}
    Some exercise goes here.
\end{exercise}

\begin{exercise}[Exercise Title]\label{exercise:test}
    Some exercise goes here.
\end{exercise}

See \myref{exercise:test}.

\section*{Miscellaneous Environments}
\subsection*{Examples Tests}
\begin{example}
    The following are statements.
    \begin{itemize}
        \item A square is a quadrilateral.
        \item 16 is a square number.
        \item 9 is a prime.
        \item Every even number greater than 2 can be expressed as the sum of two primes.
        \item $3 + 8 = 12$.
    \end{itemize}
\end{example}

\begin{example}[Statements]\label{example:test}
    The following are statements.
    \begin{itemize}
        \item A square is a quadrilateral.
        \item 16 is a square number.
        \item 9 is a prime.
        \item Every even number greater than 2 can be expressed as the sum of two primes.
        \item $3 + 8 = 12$.
    \end{itemize}
\end{example}

See \myref{example:test}.

\subsection*{Remarks Tests}
\begin{remark}
    Blah blah blah, some remark here.
\end{remark}

\begin{remark}[A Remark]
    Blah blah blah, some remark here.
\end{remark}
